\documentclass{article}
\usepackage{amsmath} 
\usepackage{amsfonts}
\usepackage{amssymb}
\usepackage[parfill]{parskip}
\usepackage{esvect}
\usepackage[thinc]{esdiff}
\usepackage{mathtools}
\usepackage{physics}
\usepackage{amsthm}
\usepackage{xfrac}
\usepackage{enumitem}

\DeclareMathOperator{\lub}{lub}
\DeclareMathOperator{\gub}{gub}

\begin{document}

\title{MATH 411: Week 5}
\author{Jacob Lockard}
\date{26 February 2026}

\maketitle

\section*{Problem 8.4}

\begin{quote}
    \itshape
    Let $f$ be convex on $(a, b)$ and let $x, y, z$ be fixed such that
    $a < x < y < z < b$. Prove that
    \begin{align*}
        \frac{(t-x)f(y) - (t-y)f(x)}{y-x} \leq f(t) \leq \frac{(z-t)f(y)+(t-y)f(z)}{z-y}
    \end{align*}
    for all $t$ with $y < t < z$, and, similarly,
    \begin{align*}
        \frac{(z-s)f(y)-(y-s)f(z)}{z-y} \leq f(s) \leq \frac{(y-s)f(x)+(s-x)f(y)}{y-x}
    \end{align*}
    for all $s$ with $x < s < y$.
\end{quote}

Let $f$ be convex on $(a, b)$ and let $x, y, z$ be fixed such that
$a < x < y < z < b$. Let $y < t < z$. Then by definition,
\begin{align*}
    f(t) \leq \frac{f(y)(z-t) + f(z)(t-y)}{z - y} \,.
\end{align*}
Also, since $x < y < t$,
\begin{align*}
    f(y) &\leq \frac{(t-y)f(x) + (y-x)f(t)}{t-x} \\
    f(y)(t-x) &\leq (t-y)f(x) + (y-x)f(t) \\
    f(y)(t-x) - (t-y)f(x) &\leq (y-x)f(t) \\
    \frac{f(y)(t-x) - (t-y)f(x)}{y-x} &\leq f(t) \,.
\end{align*}
Now let $x < s < y$. Similarly,
\begin{align*}
    \frac{(z-s)f(y)-(y-s)f(z)}{z-y} \leq f(s) \leq \frac{(y-s)f(x)+(s-x)f(y)}{y-x} \,.
\end{align*}

\qedsymbol

\section*{Problem 8.5}

\begin{quote}
    \itshape
    Let $f$ be convex on $(a,b)$. Prove that $f$ is continuous on $(a,b)$.
\end{quote}

Let $f$ be convex on $(a,b)$. We'll show that $f$ is continuous on $(a,b)$.
Let $y \in (a,b)$ and let $\epsilon > 0$. Let $x \in (a, y)$ and $z \in (y, b)$.
Denote
\begin{align*}
    m_1 = \frac{f(y) - f(x)}{y-x}\,, \quad
    m_2 = \frac{f(z) - f(y)}{z-y} \,.
\end{align*}
$m_1$ and $m_2$ exist, since $x \neq y$ and $z \neq y$. Let
\begin{align*}
    \delta = \min \Bigl\{ \frac{\epsilon}{|m_1|}, \frac{\epsilon}{|m_2|}, z-y, y-x \Bigr\} \,,
\end{align*}
removing $\frac{\epsilon}{|m_1|}$ if $m_1 = 0$ and removing $\frac{\epsilon}{|m_2|}$
if $m_2 = 0$.
$\delta$ is positive, since $\epsilon > 0$, $z > y$, and $x < y$.
Let $t \in (a,b)$ be such that $|t-y| < \delta$. If $t = y$,
then $|f(t) - f(y)| = 0 < \epsilon$.

If $t - y > 0$, then $t-y = |t-y| < \delta \leq z-y$, so
$y < t < z$, and Problem 8.4 ensures
\begin{gather*}
    \frac{(t-x)f(y) - (t-y)f(x)}{y-x} \leq f(t) \leq \frac{(z-t)f(y)+(t-y)f(z)}{z-y} \\
    \frac{(t-y)f(y) - (t-y)f(x)}{y-x} \leq f(t) - f(y) \leq \frac{-(t-y)f(y)+(t-y)f(z)}{z-y} \\
    \frac{t-y}{y-x} (f(y) - f(x)) \leq f(t) - f(y) \leq \frac{t-y}{z-y} (f(z) - f(y)) \\
    m_1 (t-y) \leq f(t) - f(y) \leq m_2 (t-y) \,.
\end{gather*}
Now, if $m_1 = 0$, we have
\begin{align*}
    |m_1(t-y)| = 0 < \epsilon \,,
\end{align*}
and if $m_1 \neq 0$, then since
$t-y = |t-y| < \frac{\epsilon}{|m_1|}$, we have
\begin{align*}
    |m_1(t-y)| &< \epsilon \,.
\end{align*}
Likewise, if $m_2 = 0$, we have
\begin{align*}
    |m_2(t-y) = 0 < \epsilon \,,
\end{align*}
and if $m_2 \neq 0$, then since $t-y = |t-y| < \frac{\epsilon}{|m_2|}$, we know
\begin{align*}
    |m_2(t-y)| &< \epsilon \,.
\end{align*}
So in any case,
\begin{align*}
    \bigl|f(t) - f(y)\bigr| < \epsilon \,.
\end{align*}
A similar argument (which we will not reproduce here) applies with the second set of inequalities in Problem 8.4
if $t < y$.

\qedsymbol

\section*{Problem 8.6}

\begin{quote}
    \itshape
    Let $f$ be defined on $[a,b]$. Prove that if $f$ is
    differentiable at a point $x \in [a,b]$, then $f$ is
    continuous at $x$.
\end{quote}

Let $f$ be defined on $[a,b]$ and differentiable at some
$x \in [a,b]$. We'll show that $f$ is continuous at $x$.

Let $c = f'(x)$, which exists since $f$ is differentiable at $x$.
Let $\epsilon > 0$.
By definition, there's some $\gamma > 0$
such that if $0 < |t-x| < \gamma$, then
\begin{gather*}
    |\phi_x(t) - c| < 1 \,.
\end{gather*}
Let
\begin{align*}
    \delta = \min \Bigl\{ \gamma, \frac{\epsilon}{1 + |c|} \Bigr\} \,,
\end{align*}
which exists and is positive since $\gamma$, $\epsilon$, and
$1 + |c|$ are all positive.

Let $t \in (a,b)$ be such that $|t-x| < \delta$. If $t = x$, then
$\bigl|f(t) - f(x)\bigr| = 0 < \epsilon$.
If $t \neq x$, then, by definition,
\begin{align*}
    f(t) - f(x) &= \phi_x(t)(t-x) \\
    &= \bigl(\phi_x(t) - c +c \bigr)(t-x) \\
    \bigl|f(t) - f(x)\bigr| &= \bigl|\bigl(\phi_x(t) - c +c \bigr)(t-x)\bigr| \\
    &= |\phi_x(t) - c +c||t-x| \\
    &\leq |t-x| \bigl( |\phi_x(t) - c| + |c| \bigr) \\
    &< |t-x| \bigl( 1 + |c| \bigr) \\
    &< \frac{\epsilon}{1 + |c|} \bigl(1 + |c|\bigr) \\
    &= \epsilon \,,
\end{align*}
where the first inequality is justified by the triangle inequality,
and the others are justified because $|t-x| < \delta$.
We conclude that $f$ is continuous at $x$.

\qedsymbol

\section*{Problem 8.7}

\begin{quote}
    \itshape
    Suppose that $f$ and $g$ are functions defined on $[a,b]$ and are both
    differentiable at a point $x \in [a,b]$. Show that $(f+g)'(x) = f'(x) + g'(x)$.
\end{quote}

Let $f$ and $g$ be defined on $[a,b]$ and differentiable at some $x \in [a,b]$.
We'll show that $(f+g)'(x) = f'(x) + g'(x)$. By definition,
\begin{align*}
    (f+g)'(x) &= \lim_{t\to x} \frac{(f+g)(t) - (f+g)(x)}{t-x} \\
    &= \lim_{t\to x} \frac{f(t) + g(t) - f(x) - g(x)}{t-x} \\
    &= \lim_{t\to x} \frac{f(t) - f(x) + g(t) - g(x)}{t-x} \\
    &= \lim_{t\to x} \Biggl( \frac{f(t) - f(x)}{t-x} + \frac{g(t) - g(x)}{t-x} \Biggr) \\
    &= \lim_{t\to x} \frac{f(t) - f(x)}{t-x} + \lim_{t\to x} \frac{g(t) - g(x)}{t-x} \\
    &= f'(x) + g'(x) \,,
\end{align*}
which exists since $f$ and $g$ are differentiable at $x$.

\qedsymbol

\section*{Problem 8.8}

\begin{quote}
    \itshape
    Suppose that $f$ and $g$ are functions defined on $[a,b]$ and are both
    differentiable at a point $x \in [a,b]$. Show that $(fg)'(x) = f'(x)g(x) + f(x)g'(x)$.
\end{quote}

Let $f$ and $g$ be defined on $[a,b]$ and differentiable at some $x \in [a,b]$.
We'll show that $(fg)'(x) = f'(x)g(x) + f(x)g'(x)$.

For all $t \in [a,b]$, we denote
\begin{align*}
    F_t = f(t) - f(x) \,; \\
    G_t = g(t) - g(x) \,.
\end{align*}

We observe that
\begin{align*}
    f(t)g(t) &= \bigl(f(x) + f(t) - f(x)\bigr) \bigl(g(x) + g(t) - g(x)\bigr) \\
    &= (f(x) + F_t)(g(x) + G_t) \\
    &= f(x)g(x) + f(x) G_t + g(x) F_t + F_t G_t \,.
\end{align*}
So,
\begin{align*}
    \lim_{t \to x} \frac{f(t)g(t) - f(x)g(x)}{t-x} &= \lim_{t \to x} \frac{f(x)G_t + g(x) F_t + F_t G_t}{t-x} \\
    &= f(x) \lim_{t \to x} \frac{G_t}{t-x} + g(x) \lim_{t \to x} \frac{F_t}{t-x} + \lim_{t \to x} F_t G_t \\
    &= f(x) g'(x) + g(x) f'(x) + \lim_{t \to x} F_t G_t \,.
\end{align*}
Problem 8.6 ensures $f$ and $g$ are continuous at $x$. So $t \mapsto F_t G_t$ is continuous
by Problem 6.6, and we can say
\begin{align*}
    \lim_{t \to x} F_t G_t
    = F_x G_x = \Bigl( f(x) - f(x) \Bigr) \Bigl( g(x) - g(x) \Bigr) = 0 \,.
\end{align*}
So we have
\begin{align*}
    (fg)'(x)
    &= \lim_{t \to x} \frac{(fg)(t) - (fg)(x)}{t-x} \\
    &= \lim_{t \to x} \frac{f(t)g(t) - f(x)g(x)}{t-x} \\
    &= f(x) g'(x) + g(x) f'(x) \,.
\end{align*}

\qedsymbol


\end{document}